% --- IMPORTS ---
\documentclass[leqno]{article}
\usepackage{graphicx}
\usepackage{dsfont}
\usepackage{mathtools}
\usepackage{amssymb}
\usepackage{amsmath}
\usepackage{amsthm}
\usepackage{float}
\usepackage{ragged2e}
\usepackage{tensor}
\usepackage{fancyhdr}
\usepackage{empheq}
\usepackage[most]{tcolorbox}
\usepackage{enumitem}
\usepackage{dirtytalk}
\usepackage{marginnote}
\usepackage{microtype}
\usepackage[
  a4paper,
  left=0.8in,
  right=1in,
  top=1in,
  bottom=0.5in,
  headheight=14pt,
  headsep=0.4in
]{geometry}
\setlength{\parindent}{0cm}
% --- TikZ Packages ---
\usepackage{pgfplots}
\pgfplotsset{compat=1.18}
\usepgfplotslibrary{fillbetween, polar}
\usetikzlibrary{calc, positioning}
\usetikzlibrary{angles, quotes}
\usetikzlibrary{arrows.meta}
\usetikzlibrary{decorations.pathreplacing, decorations.markings}
\usetikzlibrary{patterns}
\usetikzlibrary{intersections}
% --- URL Packages ---
\usepackage[colorlinks=true,linkcolor=red,citecolor=magenta,urlcolor=black]{hyperref}%
\urlstyle{same}
% --- END OF IMPORTS ---

% --- CUSTOM COMMANDS ---
\RenewDocumentCommand{\marks}{o m}{%
  \IfValueTF{#1}%
    {\marginnote{\raggedright\textbf{[#2]}}[#1]}%
    {\marginnote{\raggedright\textbf{[#2]}}}%
}
\newcommand{\vspacerule}[2]{%
  \par\vspace*{#1}%
  \noindent\hrulefill\par
  \vspace*{#2}%
}
\newcommand{\questionitem}{%
  \item\phantomsection
  \addcontentsline{toc}{section}{Question \number\value{enumi}}%
}
\renewcommand{\contentsname}{Questions}
% --- END OF CUSTOM COMMANDS ---

% --- HEADER CONFIG ---
\pagestyle{fancy}
\fancyhead{}
\fancyfoot{}
\renewcommand{\headrulewidth}{0.9pt}
\lhead{1st Order Differential Equations - Integrating Factor}
\chead{}
\rhead{\href{https://danieljsmith.org}{danieljsmith.org}}
% --- END OF HEADER CONFIG ---

\begin{document}
\vspace*{20pt}
\tableofcontents
\newpage
\begin{enumerate}[
  label=\textbf{\arabic*.},
  leftmargin=*,
  topsep=0pt,
  partopsep=0pt,
  parsep=0pt
]

% ---- Question 1 ----
\questionitem
% [Source: _QBT__Integrating_Factor_Method, Q1]

\begin{enumerate}[label=\textbf{(\alph*)}]
  \item Show that an appropriate integrating factor for
  \[
  x^2\,\frac{\mathrm{d}y}{\mathrm{d}x}+y=1,\qquad x>0
  \]
  is $e^{-1/x}$ \marks{4} \\
  \item Hence find the solution of the differential equation
  \[
  x^2\,\frac{\mathrm{d}y}{\mathrm{d}x}+y=1,\qquad x>0
  \]
  for which $y=3$ when $x=1$ \marks{6} \\
\end{enumerate}

\newpage

% ---- Question 2 ----
\questionitem
% [Source: _QBT__Integrating_Factor_Method, Q2]

\begin{enumerate}[label=\textbf{(\alph*)}]
  \item Show that
  \[
  \frac{\mathrm{d}}{\mathrm{d}x}\bigl(\ln(\cosh x+1)\bigr)=\tanh\!\left(\frac{x}{2}\right)
  \]
  \marks[-20pt]{3}
  \item Hence solve the differential equation
  \[
  \frac{\mathrm{d}y}{\mathrm{d}x}+\tanh\!\left(\frac{x}{2}\right)y=\sinh x
  \]
  for which $y=2$ when $x=0$. Give your answer in exact form \marks{7} \\
\end{enumerate}

\newpage

% ---- Question 3 ----
\questionitem
% [Source: _QBT__Integrating_Factor_Method, Q3]

\begin{enumerate}[label=\textbf{(\alph*)}]
  \item For $x>-1$, determine the general solution of the differential equation
  \[
  (1+x)\,\frac{\mathrm{d}y}{\mathrm{d}x}+y=\ln(1+x)
  \]
  giving your answer in the form $y=f(x)$ \marks{4} \\
  \item Given that $y=-1$ when $x=0$, determine the positive value of $x$ for which $y=0$ \marks{2} \\
\end{enumerate}

\newpage

% ---- Question 4 ----
\questionitem
% [Source: _QBT__Integrating_Factor_Method, Q4]

\begin{enumerate}[label=\textbf{(\alph*)}]
  \item Show that, for $0<x<\frac{\pi}{2}$, an appropriate integrating factor for
  \[
  \cos x\,\frac{\mathrm{d}y}{\mathrm{d}x}+y\sin x=x
  \]
  is $\sec x$ \marks{4} \\
  \item Hence find the solution of
  \[
  \cos x\,\frac{\mathrm{d}y}{\mathrm{d}x}+y\sin x=x
  \]
  for $0<x<\frac{\pi}{2}$, given that $y=1$ when $x=\frac{\pi}{4}$ \marks{6} \\
\end{enumerate}

\newpage

% ---- Question 5 ----
\questionitem
% [Source: _QBT__Integrating_Factor_Method, Q5]

\begin{enumerate}[label=\textbf{(\alph*)}]
  \item Show that, for $x>0$,
  \[
  \frac{\mathrm{d}}{\mathrm{d}x}\!\left(\ln\!\left(\coth \frac{x}{2}\right)\right)=-\operatorname{cosech} x
  \]
  \marks[-20pt]{3}
  \item Hence solve the differential equation
  \[
  \frac{\mathrm{d}y}{\mathrm{d}x}-\operatorname{cosech} x\,y=\tanh \frac{x}{2},\qquad x>0
  \]
  for which $y=1$ when $x=\ln 3$. Give your answer in exact form \marks{7} \\
\end{enumerate}

\newpage

% ---- Question 6 ----
\questionitem
% [Source: _QBT__Integrating_Factor_Method, Q6]

\begin{enumerate}[label=\textbf{(\alph*)}]
  \item Show that an appropriate integrating factor for
  \[
  x(x-1)\,\frac{\mathrm{d}y}{\mathrm{d}x}+(x+1)y=x(x-1)^2
  \]
  is $\dfrac{(x-1)^2}{x}$ \marks{4} \\
  \item Hence find the solution of the differential equation
  \[
  x(x-1)\,\frac{\mathrm{d}y}{\mathrm{d}x}+(x+1)y=x(x-1)^2
  \]
  for $x>1$, given that $y=2$ when $x=2$ \marks{6} \\
\end{enumerate}

\newpage

% ---- Question 7 ----
\questionitem
% [Source: _QBT__Integrating_Factor_Method, Q7]

Find the general solution of the differential equation\\
\[
\frac{\mathrm{d}y}{\mathrm{d}x}+\cot x\,y=\frac{\cos x}{\sin^2 x\left(1+\ln^2(\sin x)\right)}
\]\\
where $0<x<\pi$ \marks{7}

\newpage

% ---- Question 8 ----
\questionitem
% [Source: _QBT__Integrating_Factor_Method, Q8]

\begin{enumerate}[label=\textbf{(\alph*)}]
  \item Solve the differential equation
  \[
  (1+x)\,\frac{\mathrm{d}y}{\mathrm{d}x}-2y=(1+x)^4,\qquad x>-1
  \]
  giving your answer in the form $y=f(x)$ \marks{3} \\
  \item The solution curve passes through the point $(0,-1)$. Find the exact positive value of $x$ for which $y=0$ \marks{3} \\
\end{enumerate}

\newpage

% ---- Question 9 ----
\questionitem
% [Source: _QBT__Integrating_Factor_Method, Q9]

Find the general solution of the differential equation\\
\[
\frac{\mathrm{d}y}{\mathrm{d}x}-2y\tan x=\sec^3 x
\]\\
where $-\frac{\pi}{2}<x<\frac{\pi}{2}$ \marks{7}

\newpage

% ---- Question 10 ----
\questionitem
% [Source: _QBT__Integrating_Factor_Method, Q10]

\begin{enumerate}[label=\textbf{(\alph*)}]
  \item Show that
  \[
  \frac{\mathrm{d}}{\mathrm{d}x}\!\left(\ln\!\left(\frac{\cosh x+1}{\sinh x}\right)\right)=-\operatorname{cosech} x
  \]
  \marks[-20pt]{3}
  \item Hence solve, for $x>0$,
  \[
  \sinh x\,\frac{\mathrm{d}y}{\mathrm{d}x}-y=1-\cosh x
  \]
  given that $y=2$ when $x=\ln 3$. Give your answer in exact form \marks{7} \\
\end{enumerate}

\newpage

% ---- Question 11 ----
\questionitem
% [Source: _QBT__Integrating_Factor_Method, Q11]

\begin{enumerate}[label=\textbf{(\alph*)}]
  \item Find the general solution of the differential equation\\
  \[
  \frac{\mathrm{d}y}{\mathrm{d}x}+\frac{2x}{x^2+1}y=6x
  \]
  giving your answer in the form $y=f(x)$ \marks{3} \\
  \item Given that $y=\tfrac{7}{2}$ when $x=0$, find the smallest positive value of $x$ for which $y=4$ \marks{3} \\
\end{enumerate}

\newpage

% ---- Question 12 ----
\questionitem
% [Source: _QBT__Integrating_Factor_Method, Q12]

For $x>0$, solve the differential equation\\
\[
\sinh x\,\frac{\mathrm{d}y}{\mathrm{d}x} - 2y\cosh x = \sinh^2 x\cosh x
\] \\
given that $y=5$ when $x=\ln(1+\sqrt{2})$. \\
Give your answer in exact form. \marks{7}

\newpage

% ---- Question 13 ----
\questionitem
% [Source: _QBT__Integrating_Factor_Method, Q13]

For $x>0$, a function $y$ satisfies \\
\[
\sinh x\,\frac{\mathrm{d}y}{\mathrm{d}x}-y\cosh x=\sinh x\,\cosh^3 x
\] \\
\begin{enumerate}[label=\textbf{(\alph*)}]
  \item By first writing the equation in the form $\frac{\mathrm{d}y}{\mathrm{d}x}+P(x)y=Q(x)$, find an integrating factor \marks{2} \\
  \item Hence solve the differential equation, given that $y=2$ when $x=\ln(1+\sqrt{2})$ \marks{5} \\
\end{enumerate}
Give your answer in exact form.

\newpage

% ---- Question 14 ----
\questionitem
% [Source: _QBT__Integrating_Factor_Method, Q14]

For $x>0$, solve the differential equation\\
\[
\sinh x\,\frac{\mathrm{d}y}{\mathrm{d}x} - y\cosh x = \sinh x\cosh x
\] \\
given that $y = 2$ when $x = \ln(1+\sqrt{2})$. \\
Give your answer in exact form. \marks{7}

\newpage

% ---- Question 15 ----
\questionitem
% [Source: _QBT__Integrating_Factor_Method, Q15]

Find the general solution of the differential equation\\
\[
\frac{\mathrm{d}y}{\mathrm{d}x} - y\tan x = \frac{1}{\sqrt{1+\sin x}}
\] \\
where $-\tfrac{\pi}{2} < x < \tfrac{\pi}{2}$ \marks{7}
\end{enumerate}
\end{document}
